\begin{table}[t]
% These quantitative results indicate that competitive methods outperform ours, yet qualitative results, presented in Fig.~\ref{fig:qualiative-seg} do not seem to reflect this performance gap.
\begin{center}
\begin{small}
\begin{sc}
\begin{tabular}{cccc}
\toprule
Method & mPA(\%) & PA(\%) & IoU(\%) \\
\midrule
AttSeg & - & 47.9 & -\\
GlAtEx & - & 52.4 & -\\
Ours (MAE-weights) & 32.2 & \textbf{70.27} & 24.4\\
Ours (SETR-weights) & \textbf{35.6} & 69.5 & \textbf{27.6}\\
\bottomrule
\end{tabular}
\end{sc}
\end{small}
\end{center}
\vskip -0.1in
\caption{\textbf{Segmentation results:} Comparison of our model against AttSeg~\protect\cite{seifi2020segment} and GlAtEx~\protect\cite{seifi2021glimpse}. The metric used is mean Pixel-wise Accuracy (mPA, higher is better), Pixel-Accuracy (PA, higher is better), and Intersection over Union (IoU, higher is better). Our solution outperforms competitive methods on all metrics.}
\label{tab:segmentation}
\end{table}

              
              
              
              
              
              
              
              
              
              
              
              
              
              
              
              
          

              
              
              
              
              
              
              
              
              
              
              
              
              
              
              
              
          

              
              
              
              
              
              
              
              
              
              
              
              
              
              
              
              
          